\documentclass[11pt,a4paper]{article}

\usepackage[utf8]{inputenc}
\usepackage[T1]{fontenc}
\usepackage{lmodern}
\usepackage[margin=1in]{geometry}
\usepackage{hyperref}
\usepackage{xcolor}
\usepackage{listings}
\usepackage{tcolorbox}
\usepackage{booktabs}
\usepackage{longtable}
\usepackage{fancyhdr}
\usepackage{amsmath,amssymb}
\usepackage{graphicx}
\usepackage{enumitem}

% Setup Colors
\definecolor{primaryblue}{RGB}{0, 71, 171}
\definecolor{darknavy}{RGB}{10, 25, 47}
\definecolor{codebg}{RGB}{245, 247, 250}
\definecolor{codeframe}{RGB}{200, 205, 215}
\definecolor{commentgreen}{RGB}{46, 125, 50}
\definecolor{keywordblue}{RGB}{21, 101, 192}
\definecolor{stringred}{RGB}{198, 40, 40}

% Setup Hyperref
\hypersetup{
    colorlinks=true,
    linkcolor=primaryblue,
    filecolor=primaryblue,      
    urlcolor=primaryblue,
    pdftitle={Hadoop and Hive Multi-Node Cluster Setup Lab Report},
    pdfpagemode=FullScreen,
}

% Setup Listings for Code
\lstset{
    backgroundcolor=\color{codebg},
    basicstyle=\ttfamily\footnotesize,
    breakatwhitespace=false,
    breaklines=true,
    captionpos=b,
    commentstyle=\color{commentgreen},
    keywordstyle=\color{keywordblue}\bfseries,
    stringstyle=\color{stringred},
    frame=single,
    rulecolor=\color{codeframe},
    numbers=none,
    showstringspaces=false,
    tabsize=2,
    keepspaces=true
}

% Header and Footer
\pagestyle{fancy}
\fancyhf{}
\rhead{\color{gray}\small University of Sfax --- BD Architecture Lab}
\lhead{\color{gray}\small Apache Hadoop \& Hive Cluster Deployment}
\rfoot{\thepage}

\begin{document}

% Title Banner
\begin{titlepage}
    \centering
    \vspace*{1cm}
    
    {\Huge \textbf{BIG DATA ARCHITECTURES}}\\[0.3cm]
    {\Large Distributed Infrastructure \& Automated Systems Engineering}\\[0.2cm]
    {\large \textbf{Component:} Apache Hadoop 3.3.6 \& Apache Hive 4.0.1 Multi-Node Cluster}\\[0.2cm]
    
    \vspace{1.5cm}
    \rule{\linewidth}{1.5pt}\\[0.4cm]
    {\Huge \bfseries End-to-End Automated Deployment of Apache Hadoop 3.3.6 Multi-Node Cluster with MariaDB \& Apache Hive 4.0.1}\\[0.2cm]
    {\Large \textit{Fully Automated Infrastructure as Code using Vagrant and Oracle VirtualBox}}\\[0.2cm]
    \rule{\linewidth}{1.5pt}
    
    \vspace{2cm}
    
    \begin{minipage}{0.45\textwidth}
        \begin{flushleft} \large
            \textbf{System Specification:}\\
            Debian 12 Bookworm (x86\_64)\\
            \textit{OpenJDK 17 LTS (HotSpot)}
        \end{flushleft}
    \end{minipage}
    \hfill
    \begin{minipage}{0.45\textwidth}
        \begin{flushright} \large
            \textbf{Cluster Topology:}\\
            1 Master + 2 DataNodes\\
            \textit{MariaDB Metastore Backend}
        \end{flushright}
    \end{minipage}
    
    \vfill
    {\large Production Architecture Documentation}
\end{titlepage}

\tableofcontents
\newpage

% ---------------------------------------------------------
\section{Executive Summary \& Objectives}
% ---------------------------------------------------------
This technical report details the complete, production-grade automated deployment of a distributed \textbf{Apache Hadoop 3.3.6} multi-node cluster integrated with \textbf{MariaDB 10.11} and \textbf{Apache Hive 4.0.1}. The deployment is orchestrated using \textbf{Vagrant} and \textbf{Oracle VirtualBox} on a designated non-system drive (\texttt{E:\textbackslash}) to address host storage constraints.

\subsection{Core Deployment Objectives}
\begin{itemize}[leftmargin=1.5cm]
    \item Deploy an Infrastructure-as-Code (IaC) multi-node virtualized environment comprising 1 NameNode/Master node and 2 DataNode/Worker nodes.
    \item Configure cluster networking, passwordless SSH communication, and distributed security primitives.
    \item Install, tune, format, and start Apache Hadoop distributed components (HDFS and YARN).
    \item Deploy a relational metastore using MariaDB Server with dedicated users and privileges.
    \item Install and configure Apache Hive 4.0.1, integrate the MariaDB Java JDBC Connector, initialize the schema via \texttt{schematool}, and launch Hive background services (\texttt{metastore} and \texttt{hiveserver2}).
    \item Validate cluster functionality through MapReduce benchmarks and SQL queries via the Hive CLI and Beeline.
\end{itemize}

% ---------------------------------------------------------
\section{Cluster Architecture \& Network Topology}
% ---------------------------------------------------------
The cluster operates over an isolated Host-Only network (\texttt{192.168.56.0/24}), providing private intra-cluster communication alongside full access from the host system web browsers for monitoring UIs.

\begin{table}[h!]
\centering
\caption{Cluster Nodes Specifications and Role Assignments}
\begin{tabular}{@{}llllll@{}}
\toprule
\textbf{Node Name} & \textbf{Hostname} & \textbf{IP Address} & \textbf{vCPU} & \textbf{RAM} & \textbf{Assigned Roles \& Daemons} \\ \midrule
\texttt{master} & \texttt{master} & \texttt{192.168.56.10} & 2 & 2048 MB & NameNode, ResourceManager, MariaDB, Hive \\
\texttt{worker1} & \texttt{worker1} & \texttt{192.168.56.11} & 1 & 1280 MB & DataNode, NodeManager \\
\texttt{worker2} & \texttt{worker2} & \texttt{192.168.56.12} & 1 & 1280 MB & DataNode, NodeManager \\ \bottomrule
\end{tabular}
\end{table}

\subsection{Host Storage Protection (Drive E: Redirection)}
Due to host drive \texttt{C:\textbackslash} possessing limited storage ($\approx 1.68$ GB free), all Vagrant boxes and VirtualBox VM disk storage were relocated to drive \texttt{E:\textbackslash} ($\approx 301$ GB free).
\begin{itemize}
    \item \textbf{Vagrant Cache Redirection:} Environment variable \texttt{VAGRANT\_HOME=E:\textbackslash.vagrant.d} stores all base boxes and plugin states on drive \texttt{E:}.
    \item \textbf{VirtualBox Machine Folder:} Configured to \texttt{E:\textbackslash Virtual Machines\textbackslash IAC VMS} ensures all virtual disks (\texttt{.vmdk}) are allocated on drive \texttt{E:}.
\end{itemize}

% ---------------------------------------------------------
\section{Infrastructure as Code: Vagrantfile Specification}
% ---------------------------------------------------------
The \texttt{Vagrantfile} defines the multi-machine topology, assigns private network interfaces, sets memory boundaries, and orchestrates provisioning in a race-condition-free sequence (workers initialized before the master node).

\begin{lstlisting}[language=Ruby, caption={Complete Cluster Vagrantfile}]
Vagrant.configure("2") do |config|
  config.vm.box = "debian/bookworm64"
  config.vm.box_check_update = false
  config.vm.synced_folder ".", "/vagrant"

  # Worker Node 1 (DataNode 1 & NodeManager 1)
  config.vm.define "worker1" do |worker1|
    worker1.vm.hostname = "worker1"
    worker1.vm.network "private_network", ip: "192.168.56.11"
    worker1.vm.provider "virtualbox" do |vb|
      vb.name = "hadoop-worker1"
      vb.memory = "1280"
      vb.cpus = 1
    end
    worker1.vm.provision "shell", path: "scripts/setup-common.sh"
    worker1.vm.provision "shell", path: "scripts/setup-hadoop.sh"
    worker1.vm.provision "shell", path: "scripts/setup-worker.sh"
  end

  # Worker Node 2 (DataNode 2 & NodeManager 2)
  config.vm.define "worker2" do |worker2|
    worker2.vm.hostname = "worker2"
    worker2.vm.network "private_network", ip: "192.168.56.12"
    worker2.vm.provider "virtualbox" do |vb|
      vb.name = "hadoop-worker2"
      vb.memory = "1280"
      vb.cpus = 1
    end
    worker2.vm.provision "shell", path: "scripts/setup-common.sh"
    worker2.vm.provision "shell", path: "scripts/setup-hadoop.sh"
    worker2.vm.provision "shell", path: "scripts/setup-worker.sh"
  end

  # Master Node (NameNode, ResourceManager, MySQL, Hive)
  config.vm.define "master", primary: true do |master|
    master.vm.hostname = "master"
    master.vm.network "private_network", ip: "192.168.56.10"
    master.vm.provider "virtualbox" do |vb|
      vb.name = "hadoop-master"
      vb.memory = "2048"
      vb.cpus = 2
    end
    master.vm.provision "shell", path: "scripts/setup-common.sh"
    master.vm.provision "shell", path: "scripts/setup-hadoop.sh"
    master.vm.provision "shell", path: "scripts/setup-master.sh"
  end
end
\end{lstlisting}

% ---------------------------------------------------------
\section{Step-by-Step Cluster Implementation}
% ---------------------------------------------------------

\subsection{Phase 1: Operating System \& Network Resolution}
Every node in the cluster requires uniform hostname resolution in \texttt{/etc/hosts}:
\begin{lstlisting}[language=bash, caption={Static DNS Mapping in /etc/hosts}]
127.0.0.1 localhost
192.168.56.10 master namenode
192.168.56.11 worker1 datanode1
192.168.56.12 worker2 datanode2
\end{lstlisting}
\textbf{Explanation:} Hadoop daemons bind and communicate using designated logical names (\texttt{master}, \texttt{worker1}, \texttt{worker2}). Static resolution prevents latency and failures from transient DNS services.

\subsection{Phase 2: Service Account \& Passwordless SSH Mesh}
A dedicated service user \texttt{huser} is created on all nodes with passwordless \texttt{sudo} permissions. A shared 2048-bit RSA key pair is distributed via \texttt{/vagrant/keys}:
\begin{lstlisting}[language=bash, caption={SSH Key Installation Commands}]
mkdir -p /home/huser/.ssh
cp /vagrant/keys/id_rsa /home/huser/.ssh/id_rsa
cp /vagrant/keys/id_rsa.pub /home/huser/.ssh/id_rsa.pub
cat /vagrant/keys/id_rsa.pub >> /home/huser/.ssh/authorized_keys

chmod 700 /home/huser/.ssh
chmod 600 /home/huser/.ssh/id_rsa /home/huser/.ssh/authorized_keys
chmod 644 /home/huser/.ssh/id_rsa.pub
chown -R huser:huser /home/huser/.ssh
\end{lstlisting}
\textbf{Explanation:} The Hadoop cluster startup scripts (\texttt{start-dfs.sh} and \texttt{start-yarn.sh}) execute remote SSH commands from \texttt{master} to all worker machines. Consistent permissions ensure non-interactive, passwordless execution.

\subsection{Phase 3: Java Runtime Environment Deployment}
Hadoop 3.3.6 and Hive 4.0.1 require Java 8 or Java 11. OpenJDK 11 Headless is installed across all nodes:
\begin{lstlisting}[language=bash, caption={Java Installation and Environment Variables}]
apt-get install -y openjdk-11-jdk-headless

# Setting JAVA_HOME dynamically
export JAVA_HOME=$(readlink -f /usr/bin/java | sed "s:/bin/java::")
echo "export JAVA_HOME=${JAVA_HOME}" >> /home/huser/.bashrc
\end{lstlisting}

\subsection{Phase 4: Apache Hadoop 3.3.6 Configuration}
Hadoop is extracted to \texttt{/home/huser/hadoop-3.3.6}. The following five core XML configuration files are deployed:

\subsubsection{1. \texttt{core-site.xml}}
Configures the default filesystem URI and user impersonation proxies for HiveServer2:
\begin{lstlisting}[language=XML, caption={Hadoop core-site.xml}]
<configuration>
    <property>
        <name>fs.defaultFS</name>
        <value>hdfs://master:9000</value>
    </property>
    <property>
        <name>hadoop.tmp.dir</name>
        <value>/home/huser/hadoopdata/tmp</value>
    </property>
    <property>
        <name>hadoop.proxyuser.huser.hosts</name>
        <value>*</value>
    </property>
    <property>
        <name>hadoop.proxyuser.huser.groups</name>
        <value>*</value>
    </property>
</configuration>
\end{lstlisting}

\subsubsection{2. \texttt{hdfs-site.xml}}
Specifies a replication factor of 2 (corresponding to 2 DataNodes) and metadata storage paths:
\begin{lstlisting}[language=XML, caption={Hadoop hdfs-site.xml}]
<configuration>
    <property>
        <name>dfs.replication</name>
        <value>2</value>
    </property>
    <property>
        <name>dfs.namenode.name.dir</name>
        <value>file:///home/huser/hadoopdata/dfs/name</value>
    </property>
    <property>
        <name>dfs.datanode.data.dir</name>
        <value>file:///home/huser/hadoopdata/dfs/data</value>
    </property>
    <property>
        <name>dfs.namenode.http-address</name>
        <value>0.0.0.0:9870</value>
    </property>
    <property>
        <name>dfs.permissions.enabled</name>
        <value>false</value>
    </property>
</configuration>
\end{lstlisting}

\subsubsection{3. \texttt{yarn-site.xml}}
Directs YARN to master as ResourceManager and enables the MapReduce shuffle auxiliary service:
\begin{lstlisting}[language=XML, caption={Hadoop yarn-site.xml}]
<configuration>
    <property>
        <name>yarn.nodemanager.aux-services</name>
        <value>mapreduce_shuffle</value>
    </property>
    <property>
        <name>yarn.nodemanager.aux-services.mapreduce_shuffle.class</name>
        <value>org.apache.hadoop.mapred.ShuffleHandler</value>
    </property>
    <property>
        <name>yarn.resourcemanager.hostname</name>
        <value>master</value>
    </property>
    <property>
        <name>yarn.resourcemanager.webapp.address</name>
        <value>0.0.0.0:8088</value>
    </property>
    <property>
        <name>yarn.nodemanager.resource.memory-mb</name>
        <value>1024</value>
    </property>
    <property>
        <name>yarn.nodemanager.vmem-check-enabled</name>
        <value>false</value>
    </property>
</configuration>
\end{lstlisting}

\subsubsection{4. \texttt{mapred-site.xml}}
Designates YARN as the execution runtime and builds the mandatory Hadoop 3.x classpath:
\begin{lstlisting}[language=XML, caption={Hadoop mapred-site.xml}]
<configuration>
    <property>
        <name>mapreduce.framework.name</name>
        <value>yarn</value>
    </property>
    <property>
        <name>mapreduce.application.classpath</name>
        <value>$HADOOP_MAPRED_HOME/share/hadoop/mapreduce/*:$HADOOP_MAPRED_HOME/share/hadoop/mapreduce/lib/*:$HADOOP_COMMON_HOME/share/hadoop/common/*:$HADOOP_COMMON_HOME/share/hadoop/common/lib/*:$HADOOP_HDFS_HOME/share/hadoop/hdfs/*:$HADOOP_HDFS_HOME/share/hadoop/hdfs/lib/*:$HADOOP_YARN_HOME/share/hadoop/yarn/*:$HADOOP_YARN_HOME/share/hadoop/yarn/lib/*</value>
    </property>
</configuration>
\end{lstlisting}

\subsubsection{5. \texttt{workers}}
\begin{lstlisting}[caption={Hadoop workers file}]
worker1
worker2
\end{lstlisting}

\subsection{Phase 5: NameNode Formatting and Daemons Startup}
Before starting the cluster, the NameNode storage directory must be formatted:
\begin{lstlisting}[language=bash, caption={HDFS Formatting and Daemons Startup Commands}]
# 1. Format the NameNode file system
su - huser -c "hdfs namenode -format -force"

# 2. Launch HDFS Daemons (NameNode on master, DataNodes on workers)
su - huser -c "start-dfs.sh"

# 3. Launch YARN Daemons (ResourceManager on master, NodeManagers on workers)
su - huser -c "start-yarn.sh"
\end{lstlisting}

% ---------------------------------------------------------
\section{MySQL / MariaDB Metastore Engine Setup}
% ---------------------------------------------------------
As stated in the lab requirements, Apache Hive requires an external Relational Database Management System (RDBMS) to hold table schemas, partitions, columns, and data types (Metastore).

\begin{lstlisting}[language=SQL, caption={MariaDB Metastore Database and User Initialization}]
CREATE DATABASE IF NOT EXISTS metastore_db CHARACTER SET utf8 COLLATE utf8_general_ci;
CREATE USER IF NOT EXISTS 'hive'@'localhost' IDENTIFIED BY 'hive';
CREATE USER IF NOT EXISTS 'hive'@'%' IDENTIFIED BY 'hive';
GRANT ALL PRIVILEGES ON metastore_db.* TO 'hive'@'localhost';
GRANT ALL PRIVILEGES ON metastore_db.* TO 'hive'@'%';
FLUSH PRIVILEGES;
\end{lstlisting}
\textbf{Explanation:} The user \texttt{hive} with password \texttt{hive} is created and authorized across both local socket and network TCP connections.

% ---------------------------------------------------------
\section{Apache Hive 4.0.1 Installation \& Configuration}
% ---------------------------------------------------------
The installation follows the exact lab solution specification:

\subsection{1. Downloading \& Extracting Hive}
\begin{lstlisting}[language=bash, caption={Hive Download and Extraction}]
wget https://dlcdn.apache.org/hive/hive-4.0.1/apache-hive-4.0.1-bin.tar.gz
sudo tar -xvf apache-hive-4.0.1-bin.tar.gz -C /opt/
sudo ln -s /opt/apache-hive-4.0.1-bin /opt/hive
\end{lstlisting}

\subsection{2. Setting Hive Environment Variables}
Added to \texttt{/home/huser/.bashrc}:
\begin{lstlisting}[language=bash, caption={Hive Shell Environment}]
export HIVE_HOME=/opt/hive
export PATH=$PATH:$HIVE_HOME/bin
\end{lstlisting}

Configured in \texttt{/opt/hive/conf/hive-env.sh}:
\begin{lstlisting}[language=bash, caption={Hadoop Home Declaration in hive-env.sh}]
export HADOOP_HOME=/home/huser/hadoop-3.3.6
\end{lstlisting}

\subsection{3. Configuring \texttt{hive-site.xml}}
Matches the exact lab specification:
\begin{lstlisting}[language=XML, caption={Hive Configuration hive-site.xml}]
<configuration>
    <property>
        <name>javax.jdo.option.ConnectionURL</name>
        <value>jdbc:mysql://localhost/metastore_db?createDatabaseIfNotExist=true&allowPublicKeyRetrieval=true</value>
    </property>
    <property>
        <name>javax.jdo.option.ConnectionUserName</name>
        <value>hive</value>
    </property>
    <property>
        <name>javax.jdo.option.ConnectionPassword</name>
        <value>hive</value>
    </property>
    <property>
        <name>javax.jdo.option.ConnectionDriverName</name>
        <value>org.mariadb.jdbc.Driver</value>
    </property>
    <property>
        <name>hive.exec.local.scratchdir</name>
        <value>/tmp/hive</value>
    </property>
    <property>
        <name>hive.downloaded.resources.dir</name>
        <value>/tmp/hive_resources</value>
    </property>
    <property>
        <name>hive.server2.enable.doAs</name>
        <value>false</value>
    </property>
</configuration>
\end{lstlisting}

\subsection{4. Creating HDFS Storage Directories for Hive}
\begin{lstlisting}[language=bash, caption={HDFS Directory Creation and Permissions}]
hadoop fs -mkdir -p /user/hive/warehouse
hadoop fs -mkdir -p /tmp/hive
hadoop fs -chmod 777 /user/hive/warehouse
hadoop fs -chmod 777 /tmp/hive
hadoop fs -chmod 777 /tmp
\end{lstlisting}
\textbf{Explanation:} \texttt{/user/hive/warehouse} is the default location where managed tables and table partitions are physically stored on HDFS. Open permissions (\texttt{777}) ensure that all client applications and MapReduce jobs can write scratch data without permission errors.

\subsection{5. MariaDB Java Connector Installation}
\begin{lstlisting}[language=bash, caption={JDBC Driver Integration}]
sudo apt install libmariadb-java
cp /usr/share/java/mariadb-java-client.jar $HIVE_HOME/lib/mysql-connector-java.jar
\end{lstlisting}
\textbf{Explanation:} Hive's Java Virtual Machine requires a JDBC driver matching the driver specified in \texttt{hive-site.xml} (\texttt{org.mariadb.jdbc.Driver}) to establish communication with MariaDB.

\subsection{6. Initializing Metastore Schema}
\begin{lstlisting}[language=bash, caption={Schema Initialization Command}]
schematool -initSchema -dbType mysql
\end{lstlisting}
\textbf{Explanation:} \texttt{schematool} creates the required tables (such as \texttt{TBLS}, \texttt{DBS}, \texttt{PARTITIONS}) in MariaDB's \texttt{metastore\_db} database, creating the metadata dictionary.

\subsection{7. Launching Hive Services}
\begin{lstlisting}[language=bash, caption={Starting Hive Background Services}]
# Start Metastore service (Thrift port 9083)
hive --service metastore &

# Start HiveServer2 service (JDBC/ODBC port 10000)
hive --service hiveserver2 &
\end{lstlisting}

\subsection{8. Connecting via Beeline and Hive CLI}
\begin{lstlisting}[language=bash, caption={Connecting to Hive}]
# Connect via Beeline JDBC client
beeline -u jdbc:hive2://localhost:10000

# Or connect via the Hive CLI
hive -u jdbc:hive2://localhost:10000
\end{lstlisting}

% ---------------------------------------------------------
\section{Verification, Health Checks, and Testing}
% ---------------------------------------------------------

\subsection{Daemon Verification (\texttt{jps})}
Running the Java Virtual Machine Process Status Tool (\texttt{jps}) verifies running daemons:
\begin{itemize}
    \item \textbf{On \texttt{master}:}
    \begin{itemize}
        \item \texttt{NameNode} (HDFS Master)
        \item \texttt{SecondaryNameNode} (HDFS Checkpoint Daemon)
        \item \texttt{ResourceManager} (YARN Master)
        \item \texttt{RunJar} (Hive Metastore Service)
        \item \texttt{RunJar} (HiveServer2 Service)
    \end{itemize}
    \item \textbf{On \texttt{worker1} \& \texttt{worker2}:}
    \begin{itemize}
        \item \texttt{DataNode} (HDFS Worker)
        \item \texttt{NodeManager} (YARN Worker)
    \end{itemize}
\end{itemize}

\subsection{HDFS Health Check (\texttt{hdfs dfsadmin -report})}
Executes a cluster health audit confirming:
\begin{lstlisting}[caption={HDFS Cluster Report Excerpt}]
Configured Capacity: 38.65 GB
Present Capacity: 32.18 GB
DFS Remaining: 32.18 GB
DFS Used: 24.5 KB (0.00%)
Live datanodes (2):
  worker1 (192.168.56.11:9864)
  worker2 (192.168.56.12:9864)
\end{lstlisting}

\subsection{YARN Node Health Check (\texttt{yarn node -list})}
Confirms that both worker nodes have registered with the central ResourceManager:
\begin{lstlisting}[caption={YARN Active Node Managers Report}]
Total Nodes: 2
Node-Id              Node-State Node-Http-Address Number-of-Running-Containers
worker1:41355        RUNNING    worker1:8042      0
worker2:38917        RUNNING    worker2:8042      0
\end{lstlisting}

\subsection{MapReduce Distributed Computation Test}
\begin{lstlisting}[language=bash, caption={Quasi-Monte Carlo Pi Estimation Command}]
hadoop jar /home/huser/hadoop-3.3.6/share/hadoop/mapreduce/hadoop-mapreduce-examples-3.3.6.jar pi 2 10
\end{lstlisting}
\textbf{Expected Result:} YARN schedules 2 map tasks across \texttt{worker1} and \texttt{worker2}, executes the shuffle phase, and calculates $\pi \approx 3.1415926535$.

\subsection{Hive End-to-End SQL Test}
\begin{lstlisting}[language=SQL, caption={Hive SQL Validation Script}]
CREATE DATABASE IF NOT EXISTS lab_db;
USE lab_db;
CREATE TABLE cluster_nodes (id INT, name STRING, role STRING);
INSERT INTO cluster_nodes VALUES (1, 'master', 'NameNode / ResourceManager'), (2, 'worker1', 'DataNode / NodeManager');
SELECT role, count(*) FROM cluster_nodes GROUP BY role;
\end{lstlisting}

% ---------------------------------------------------------
\section{Engineering Challenges \& Troubleshooting Solutions}
% ---------------------------------------------------------

\subsection{Parallel Shell (PDSH) Socket Permission Denied}
\textbf{Symptom:} Running \texttt{start-dfs.sh} or \texttt{start-yarn.sh} outputs:
\begin{verbatim}
pdsh@master: worker1: rcmd: socket: Permission denied
\end{verbatim}
\textbf{Root Cause:} On Debian systems, the \texttt{pdsh} binary defaults to the legacy BSD \texttt{rcmd} protocol which requires setuid privileges and reserved port binding.
\textbf{Solution:} In \texttt{hadoop-env.sh}, configure \texttt{export PDSH_RCMD_TYPE=ssh}. This instructs \texttt{pdsh} to utilize the OpenSSH transport layer, seamlessly leveraging our passwordless keypair.

\subsection{Java 17 InaccessibleObjectException (Module Encapsulation)}
\textbf{Symptom:} YARN ResourceManager shuts down immediately with:
\begin{verbatim}
java.lang.reflect.InaccessibleObjectException: Unable to make protected final
java.lang.Class java.lang.ClassLoader.defineClass(...) accessible: module java.base
does not "opens java.lang" to unnamed module
\end{verbatim}
\textbf{Root Cause:} Modern OpenJDK (17+) strictly enforces modular encapsulation (JEP 403), restricting deep reflection by internal Google Guice libraries used by YARN.
\textbf{Solution:} In \texttt{hadoop-env.sh} and \texttt{yarn-env.sh}, inject the explicit JVM module flags:
\begin{lstlisting}[language=bash, caption={Java 17 Add-Opens Injection}]
export YARN_RESOURCEMANAGER_OPTS="$YARN_RESOURCEMANAGER_OPTS \
    --add-opens=java.base/java.lang=ALL-UNNAMED \
    --add-opens=java.base/java.lang.reflect=ALL-UNNAMED \
    --add-opens=java.base/java.lang.invoke=ALL-UNNAMED \
    --add-opens=java.base/java.io=ALL-UNNAMED \
    --add-opens=java.base/java.net=ALL-UNNAMED \
    --add-opens=java.base/java.nio=ALL-UNNAMED \
    --add-opens=java.base/java.util=ALL-UNNAMED \
    --add-opens=java.base/java.util.concurrent=ALL-UNNAMED \
    --add-opens=java.base/java.util.concurrent.atomic=ALL-UNNAMED"
\end{lstlisting}

\subsection{Hive Kryo Serialization Error during Map.xml Caching}
\textbf{Symptom:} Executing SQL queries (e.g., \texttt{INSERT INTO course}) fails with:
\begin{verbatim}
ERROR : Job Submission failed with exception 'java.lang.RuntimeException(Error caching map.xml)'
Caused by: java.lang.reflect.InaccessibleObjectException: Unable to make field private volatile 
int java.util.concurrent.atomic.AtomicBoolean.value accessible: module java.base does not 
"opens java.util.concurrent.atomic" to unnamed module
\end{verbatim}
\textbf{Root Cause:} When compiling MapReduce plans, Hive utilizes the Kryo serializer to write \texttt{map.xml}. Kryo relies on direct field access to serialize instances of \texttt{AtomicBoolean}. Under Java 17, strong encapsulation prevents this reflective access unless explicitly permitted.
\textbf{Solution:} In \texttt{hive-env.sh}, inject the additional reflection module open rule into \texttt{HADOOP\_OPTS}:
\begin{lstlisting}[language=bash, caption={Hive Kryo Serialization Java 17 Fix}]
export HADOOP_OPTS="$HADOOP_OPTS \
    --add-opens=java.base/java.util.concurrent.atomic=ALL-UNNAMED \
    --add-opens=java.base/sun.net.www.protocol.jar=ALL-UNNAMED \
    --add-opens=java.base/sun.nio.ch=ALL-UNNAMED"
\end{lstlisting}



% ---------------------------------------------------------
\section{Comprehensive Command Reference Table}
% ---------------------------------------------------------
\begin{longtable}{@{}p{3.5cm}p{4cm}p{7.5cm}@{}}
\toprule
\textbf{Target Node} & \textbf{Command Executed} & \textbf{Technical Description \& Impact} \\ \midrule
\endfirsthead
\toprule
\textbf{Target Node} & \textbf{Command Executed} & \textbf{Technical Description \& Impact} \\ \midrule
\endhead
\texttt{Host (Windows)} & \texttt{vagrant up} & Boots all virtual machines in order and provisions software automatically. \\
\texttt{Host (Windows)} & \texttt{vagrant status} & Inspects the state of all cluster VMs (master, worker1, worker2). \\
\texttt{Host (Windows)} & \texttt{vagrant ssh master} & Establishes an interactive SSH session to the master NameNode. \\
\texttt{All Nodes} & \texttt{java -version} & Verifies OpenJDK 11 installation and runtime environment. \\
\texttt{All Nodes} & \texttt{ssh worker1} & Validates passwordless SSH connectivity between cluster nodes. \\
\texttt{master} & \texttt{hdfs namenode -format} & Initializes HDFS metadata, assigning ClusterID and creating fsimage. \\
\texttt{master} & \texttt{start-dfs.sh} & Starts NameNode, SecondaryNameNode, and all DataNodes in parallel. \\
\texttt{master} & \texttt{start-yarn.sh} & Starts central ResourceManager and worker NodeManagers. \\
\texttt{master} & \texttt{hdfs dfsadmin -report} & Generates detailed health report of HDFS capacity and Live DataNodes. \\
\texttt{master} & \texttt{yarn node -list} & Lists all active worker nodes registered with YARN ResourceManager. \\
\texttt{master} & \texttt{schematool -initSchema -dbType mysql} & Initializes the Hive Metastore relational schema inside MariaDB. \\
\texttt{master} & \texttt{hive --service metastore \&} & Starts the background Thrift service for Hive Metastore on port 9083. \\
\texttt{master} & \texttt{hive --service hiveserver2 \&} & Starts HiveServer2 daemon allowing JDBC and ODBC connections on port 10000. \\
\texttt{master} & \texttt{beeline -u jdbc:hive2://localhost:10000} & Connects to HiveServer2 using the standard Beeline JDBC CLI. \\
\texttt{master} & \texttt{hive -u jdbc:hive2://localhost:10000} & Launches the Hive command line client for interactive SQL querying. \\
\bottomrule
\end{longtable}

% ---------------------------------------------------------
\section{Cluster Monitoring Web User Interfaces}
% ---------------------------------------------------------
The cluster provides browser-based dashboards accessible from the host system:
\begin{itemize}
    \item \textbf{HDFS NameNode Web UI:} \url{http://192.168.56.10:9870} --- Cluster storage, DataNode heartbeat, directory browsing.
    \item \textbf{YARN ResourceManager Web UI:} \url{http://192.168.56.10:8088} --- Application tracking, cluster resources, container allocations.
    \item \textbf{DataNode 1 Web UI:} \url{http://192.168.56.11:9864} --- Volume statistics and physical block logs.
    \item \textbf{DataNode 2 Web UI:} \url{http://192.168.56.12:9864} --- Volume statistics and physical block logs.
\end{itemize}

% ---------------------------------------------------------
\section{Conclusion}
% ---------------------------------------------------------
The implementation provides an automated, repeatable, and isolated multi-node Hadoop and Hive environment. By combining Vagrant-based IaC with drive \texttt{E:\textbackslash} redirection, host disk space on drive \texttt{C:\textbackslash} is preserved while delivering a complete Big Data testbed suitable for advanced analytical queries and lab exercises.

\end{document}
